\chapter{曲線の面積と変化をどう求めるか}\label{章:導入}
\section{直線でできることを曲線へ広げる}
長方形の面積は横の長さと縦の長さの積で求められる。直線の傾きも，横の変化量に対する縦の変化量の比で求められる。ところが，境界が曲線になったり，傾きが場所ごとに変わったりすると，同じ計算をそのまま使うことはできない。この本では，直線や長方形について知っている計算から出発し，曲線の変化と面積を調べる方法を作る。

まず，二つの問いを並べてみよう。第一は，放物線 $y=x^2$ 上の点 $(1,1)$ で，曲線はどの程度の傾きで進んでいるのかという問いである。第二は，この放物線と $x$ 軸，直線 $x=1$ に囲まれた部分の面積はいくらかという問いである。第一の問いでは二点を結ぶ直線を考え，その二点を近づける。第二の問いでは細い長方形を並べ，その幅を小さくする。どちらも，計算できるものから始めて，細かくした先にある値を求めようとしている。

ここで「限りなく近づける」という表現は，まだ計算の見通しを説明するために使っている。何をどのように近づけるのかを明らかにするのが\chapterref{章:極限と微分}であり，長方形の和を積分として定義するのが\chapterref{章:定積分}である。この章では具体的な図と計算を通して，後の定義が何のために必要なのかを確かめる。

\section{二点間の変化と一点での変化}\label{節:導入変化}
関数 $y=f(x)$ において，$x$ が $a$ から $a+h$ まで変わったとする。$h\ne0$ のとき，横の変化量は $h$，縦の変化量は $f(a+h)-f(a)$ だから，二点を結ぶ直線の傾きは
\[
 \frac{f(a+h)-f(a)}{h}
\]
である。これをその区間の\textbf{平均変化率}という。「平均」と呼ぶのは，途中の変化のしかたを区別せず，始点と終点の変化だけで一つの数を作っているからである。

$f(x)=x^2$，$a=1$ では，
\[
 \frac{(1+h)^2-1}{h}
 =\frac{2h+h^2}{h}=2+h \qquad(h\ne0)
\]
となる。$h=1$ なら傾きは $3$，$h=0.1$ なら $2.1$，$h=-0.1$ なら $1.9$ である。$h$ を正の側からも負の側からも $0$ に近づけると，傾きは $2$ に近づく。この値を，放物線上の点 $(1,1)$ における接線の傾きとして採用したい。

\begin{figure}[htbp]
\centering
\begin{tikzpicture}[x=2.6cm,y=1.5cm,>=Stealth]
 \draw[->] (-0.12,0)--(2.1,0) node[right]{$x$};
 \draw[->] (0,-0.25)--(0,3.75) node[above]{$y$};
 \draw[very thick,domain=0:1.9,samples=80] plot (\x,{\x*\x}) node[right]{$y=x^2$};
 \draw[blue,thick,domain=0.45:1.75] plot (\x,{2*\x-1});
 \draw[red,thick,domain=0.55:1.8] plot (\x,{2.6*\x-1.6});
 \fill (1,1) circle(1.6pt) node[below right]{$P(1,1)$};
 \fill (1.6,2.56) circle(1.6pt) node[right]{$Q$};
 \draw[dashed] (1,0) node[below]{$1$}--(1,1);
 \draw[dashed] (1.6,0) node[below]{$1+h$}--(1.6,2.56);
 \node[blue] at (1.82,2.08) {接線};
 \node[red] at (0.65,1.6) {割線};
\end{tikzpicture}
\caption{二点を結ぶ割線の傾きを，点 $Q$ を点 $P$ に近づけて調べる。}
\label{図:導入接線}
\end{figure}

図\ref{図:導入接線}のように，曲線上の二点を結ぶ直線を\textbf{割線}と呼ぶ。割線は二点で決まるのに対し，接線の傾きは一つの点の近くの関数の振る舞いで決まる。図だけを見て「曲線に一度だけ触れる直線」と定義すると，接線が別の場所で再び曲線と交わる場合を扱えない。微分では，割線の傾きの極限を使って接線を定める。

\begin{example}{平均の速さから瞬間の速さへ}
一直線上を動く物体の時刻 $t$ 秒における位置が $s(t)=t^2$ メートルで表されている。時刻 $2$ 秒から $2+h$ 秒までの平均の速度を求め，時刻 $2$ 秒における速度を予想せよ。
\end{example}
位置の変化を経過時間で割ればよいので，平均の速度は
\[
 \frac{s(2+h)-s(2)}{h}
 =\frac{(2+h)^2-4}{h}=4+h
\]
メートル毎秒である。たとえば $h=0.1$ なら $4.1$，$h=0.01$ なら $4.01$ となる。時刻 $2$ 秒より前の区間を調べるために $h<0$ としても，同じ式が使える。どちらから近づけても $4$ に近づくので，瞬間の速度は $4$ メートル毎秒と予想できる。

この計算では $h$ を $0$ に代入していない。元の式では $h=0$ にすると分母も分子も $0$ となり，割り算ができない。$h\ne0$ のまま式を整理し，$h$ が $0$ に近づくときの値を調べたのである。\textbf{変化量を $0$ にすることと，$0$ に近づけることは別の操作}である。

\section{長方形で面積を挟む}\label{節:導入長方形}
次に，$y=x^2$ のグラフの下で，$0\le x\le1$ の範囲にある部分の面積を考える。この部分の面積を $S$ と書く。横の区間を $n$ 等分すると，一つの小区間の幅は $\frac1n$ である。$x^2$ はこの範囲で増加するから，小区間の左端の高さで作った長方形は曲線の下に収まり，右端の高さで作った長方形は曲線の上まで覆う。

左端の高さを使った和を $L_n$，右端の高さを使った和を $U_n$ とすると，
\begin{align*}
 L_n&=\frac1n\sum_{k=0}^{n-1}\left(\frac{k}{n}\right)^2,\\
 U_n&=\frac1n\sum_{k=1}^{n}\left(\frac{k}{n}\right)^2
\end{align*}
である。添字が異なるのは，左端では $0$ から $\frac{n-1}{n}$ まで，右端では $\frac1n$ から $1$ までの高さを使うためである。図\ref{図:導入長方形}では $n=4$ としている。

\begin{figure}[htbp]
\centering
\begin{tikzpicture}[x=3.7cm,y=3.1cm,>=Stealth]
 \begin{scope}
  \foreach \k in {0,1,2,3}{
   \pgfmathsetmacro{\xl}{\k/4}
   \pgfmathsetmacro{\xr}{(\k+1)/4}
   \pgfmathsetmacro{\hh}{(\k/4)^2}
   \filldraw[fill=blue!16,draw=blue!60] (\xl,0) rectangle (\xr,\hh);
  }
  \draw[->] (-0.08,0)--(1.13,0) node[right]{$x$};
  \draw[->] (0,-0.05)--(0,1.17) node[above]{$y$};
  \draw[thick,domain=0:1,samples=60] plot (\x,{\x*\x});
  \node[below] at (1,0) {$1$};
  \node at (0.5,-0.23) {左端の高さ：$L_4$};
 \end{scope}
 \begin{scope}[xshift=5.4cm]
  \foreach \k in {0,1,2,3}{
   \pgfmathsetmacro{\xl}{\k/4}
   \pgfmathsetmacro{\xr}{(\k+1)/4}
   \pgfmathsetmacro{\hh}{((\k+1)/4)^2}
   \filldraw[fill=red!13,draw=red!60] (\xl,0) rectangle (\xr,\hh);
  }
  \draw[->] (-0.08,0)--(1.13,0) node[right]{$x$};
  \draw[->] (0,-0.05)--(0,1.17) node[above]{$y$};
  \draw[thick,domain=0:1,samples=60] plot (\x,{\x*\x});
  \node[below] at (1,0) {$1$};
  \node at (0.5,-0.23) {右端の高さ：$U_4$};
 \end{scope}
\end{tikzpicture}
\caption{下からの長方形和と上からの長方形和。分割を細かくすると両者の差が小さくなる。}
\label{図:導入長方形}
\end{figure}

正方形の和の公式
\[
 \sum_{k=1}^{n}k^2=\frac{n(n+1)(2n+1)}6
\]
を使うと，
\begin{align*}
 U_n&=\frac{n(n+1)(2n+1)}{6n^3}
     =\frac13+\frac{1}{2n}+\frac{1}{6n^2},\\
 L_n&=\frac{(n-1)n(2n-1)}{6n^3}
     =\frac13-\frac{1}{2n}+\frac{1}{6n^2}
\end{align*}
となる。面積は両者の間にあるから，
\[
 L_n\le S\le U_n
\]
である。そして $n$ を大きくすれば，$L_n$ も $U_n$ も $\frac13$ に近づく。したがって $S=\frac13$ と考えられる。この議論を，曲線に囲まれた面積の定義まで含めて明確にするのが定積分である。

二つの和の差は，展開せずに求めることもできる。同じ高さの項が打ち消し合うので，
\[
 U_n-L_n
 =\frac1n\left(1^2-0^2\right)=\frac1n
\]
となる。つまり，求めたい値を幅 $\frac1n$ の範囲に閉じ込められる。このように\textbf{値そのものをただ予想するだけでなく，どれだけの幅で挟めたかを調べる}と，近似の確かさが分かる。

\begin{example}{一回の近似から誤差の範囲を読む}
上の放物線の面積を $10$ 等分した長方形で挟め。さらに，右端の和 $U_{10}$ を面積の近似値としたとき，どの程度大きく見積もる可能性があるか説明せよ。
\end{example}
$n=10$ を代入すると，
\[
 L_{10}=\frac13-\frac1{20}+\frac1{600}
       =\frac{57}{200}=0.285,
 \qquad
 U_{10}=\frac{77}{200}=0.385
\]
である。したがって $0.285\le S\le0.385$ と分かる。$U_{10}$ は常に面積以上であり，
\[
 0\le U_{10}-S\le U_{10}-L_{10}=0.1
\]
である。この情報だけからでも，大きく見積もる誤差は $0.1$ 以下と保証できる。実際の面積 $\frac13$ を知った後なら，誤差は $\frac{31}{600}$ と計算できるが，未知の値を求める途中でも使えるのは前者の評価である。

\section{面積の増え方に現れる高さ}\label{節:導入基本定理}
ここまでの二つの問いは，一見すると別の話である。接線では縦の変化を横の変化で割り，面積では縦の長さと横の長さを掛けて足している。両者をつなぐために，右端を固定せず動かしてみる。

$y=x^2$ の下で，$0$ から $x$ までの面積を $A(x)$ と書く。ただし，この段階では $x\ge0$ とする。右端を $x$ から $x+h$ へ少し動かすと，面積の増加 $A(x+h)-A(x)$ は，幅 $h$ の細い帯の面積である。$h>0$ なら，その帯の高さは $x^2$ 以上 $(x+h)^2$ 以下だから，
\[
 x^2h\le A(x+h)-A(x)\le(x+h)^2h.
\]
正の数 $h$ で割ると，
\[
 x^2\le\frac{A(x+h)-A(x)}h\le(x+h)^2.
\]
$h$ を $0$ に近づけると左右はともに $x^2$ へ近づく。面積の平均変化率も，その間にあるので $x^2$ へ近づくはずである。つまり，\textbf{面積の瞬間の増加率は，その場所の曲線の高さになる}。

この結び付きによって，長方形の和を毎回計算せずに面積を求める道が開ける。もし $A(x)=\frac{x^3}{3}$ なら，
\begin{align*}
 \frac{A(x+h)-A(x)}h
 &=\frac{(x+h)^3-x^3}{3h}\\
 &=x^2+xh+\frac{h^2}{3}
\end{align*}
であり，確かに $h$ を $0$ へ近づけると $x^2$ になる。さらに $A(0)=0$ も満たす。この計算は面積の候補が条件に合うことを示しているが，その条件で面積が一つに決まることはまだ証明していない。後で平均値の定理を学ぶと，「微分して同じ関数になる二つの関数の差は定数である」と示せる。それを使って，この候補が面積そのものであると結論できる。

ここでは，候補を見つける計算と，候補が正しいと確かめる論理を分けておこう。積分の計算で「微分の逆を探す」という方針を使えるのは，この両方がそろうからである。

\section{符号と単位を読み落とさない}
グラフの高さが負になる場合，幅と高さの積も負になる。幾何学的な面積は負にはならないが，積分で足し合わせる量には符号が付く。この違いは，運動を考えると自然に理解できる。

右向きを正とする直線上で，物体が毎秒 $2$ メートルの速度で $3$ 秒進み，その後，毎秒 $1$ メートルの速度で逆向きに $2$ 秒進んだとする。位置の変化は
\[
 2\times3+(-1)\times2=4
\]
メートルである。一方，実際に動いた道のりは
\[
 2\times3+1\times2=8
\]
メートルとなる。速度のグラフで符号を付けて面積を足すと位置の変化が得られ，速さのグラフの面積を足すと道のりが得られる。同じグラフを使っていても，何を求めているかによって計算が変わる。

単位も計算の意味を教えてくれる。位置を時間で割れば速度になり，速度に時間を掛ければ位置の変化になる。面積を横の長さで割れば縦の長さになり，縦の長さに横の長さを掛ければ面積になる。微分や積分の式が複雑になっても，単位を確かめる習慣は役に立つ。

これから，まず変化率を求める方法を整え，続いて積分を定義する。その後の各章では，式をどのように変えれば既知の微分の逆を見つけられるか，あるいは不定積分を求めずに定積分の値を決められるかを調べる。最初からすべての公式を覚える必要はない。今扱っている式で，足しているのか，割っているのか，どの変数を変えているのかを確かめながら進めよう。

\section*{章末問題}
\begin{enumerate}[label=\textbf{\thechapter.\arabic*},leftmargin=*]
 \item\label{問:導入速度} $s(t)=t^2$ とする。時刻 $3$ から $3+h$ までの平均の速度を求めよ。$h=0.2,-0.2$ の場合の値を求め，時刻 $3$ の瞬間の速度を予想せよ。
 \item\label{問:導入直線面積} $y=2x+1$ と $x$ 軸，直線 $x=0,x=2$ に囲まれた台形の面積を求めよ。また，区間 $[0,2]$ を $n$ 等分して左端・右端の高さから作る長方形和 $L_n,U_n$ を求め，それぞれが同じ面積に近づくことを確かめよ。
 \item\label{問:導入中点} $y=x^2$ の下の $0\le x\le1$ の面積を，$4$ 等分した各小区間の中央の高さで近似せよ。得られた値を $L_4,U_4$ と比較せよ。
 \item\label{問:導入三次割線} $y=x^3$ のグラフ上の二点 $(1,1)$，$(1+h,(1+h)^3)$ を結ぶ直線の傾きを求めよ。$h\ne0$ の条件を使う箇所を明記し，$h$ を $0$ に近づけたときの傾きを予想せよ。
 \item\label{問:導入拡大面積} $y=x^2$ の下の $0\le x\le2$ の面積を $S$ とする。$4$ 等分した左端と右端の長方形和を求めて $S$ を挟め。一般の $n$ 等分では，$0\le x\le1$ の場合の長方形和の何倍になるか。そこから $S$ を予想せよ。
 \item\label{問:導入符号} 直線上を動く物体の速度が，時刻 $0$ から $4$ までは毎秒 $3$ メートル，時刻 $4$ から $7$ までは毎秒 $-2$ メートルである。位置の変化と道のりをそれぞれ求めよ。二つの答えが異なる理由を，速度の符号を使って説明せよ。
\end{enumerate}

\smallskip\noindent 解答は\pageref{解答:第1章}ページ。
